% Part: sets-functions-relations
% Chapter: infinite
% Section: dedekind-induction

\documentclass[../../../include/open-logic-section]{subfiles}

\begin{document}

\olfileid{sfr}{infinite}{induction}
\olsection[Induksi Aritmetis]{Aljabar Dedekind lan Induksi Aritmetis}

Bab kang wigati saiki, aljabar Dedekind---malah, \emph{saben} aljabar
Dedekind---bisa dadi sulih kanggo wilangan natural. Iki amarga dudutan
prasaja iki:

\begin{thm}[Induksi aritmetis]\ollabel{thm:dedinfiniteinduction}
Upamakna $N, s, o$ mbentuk aljabar Dedekind. Banjur kanggo saben
himpunan $X$:
	\begin{center}
		yen $o \in X$ lan $(\forall {n} \in N \cap X){s}({n}) \in X$, {mula} $N \subseteq X$.
	\end{center}
\end{thm}

\begin{proof}
Miturut definisi aljabar Dedekind, $N = \closureofunder{s}{o}$.
Saiki yen ${o} \in X$ lan $(\forall {n} \in N)(n \in X \lif
{s}({n}) \in X)$, mula $N = \closureofunder{s}{o} \subseteq X$.
\end{proof}

Amarga induksi iku ciri mligi wilangan natural, pokoke mangkene. Saka
saben himpunan kang tanpa wates miturut Dedekind, kita bisa mbentuk
aljabar Dedekind lan nggunakake aljabar iku minangka sulih kanggo
wilangan natural.

Pancen, \olref{thm:dedinfiniteinduction} ngrumusake induksi nganggo
istilah \emph{teori himpunan}. Nanging asas iki gampang ditulis nganggo
istilah kang bisa uga luwih kulina:

\begin{cor}\ollabel{natinductionschema}
Upamakna $N, s, o$ mbentuk aljabar Dedekind. Banjur kanggo saben formula
$\phi(x)$, kang bisa duwe paramèter:
\begin{center}
	yen $\phi(o)$ lan $(\forall {n} \in N)(\phi(n)\lif
	\phi({s}({n})))$, {mula} $(\forall n \in N)\phi(n)$
\end{center}
\end{cor}

\begin{proof}
Jupuken $X = \Setabs{n \in N}{\phi(n)}$, banjur gunakna
\olref{thm:dedinfiniteinduction}
\end{proof}

Ing asil iki, kita nyebut formula kang ``duwe paramèter''. Tegese,
kanthi kasar, kanggo saben obyek $c_1, \ldots, c_k$, kita bisa nggarap
$\phi(x, c_1, \ldots, c_k)$. Kanthi luwih pas, asil kasebut bisa
diandharake tanpa nyebut ``paramèter'' kaya mangkene. Kanggo saben
formula $\phi(x, v_1, \ldots, v_k)$, kang kabeh variabel bebane wis
dituduhake, kita duwe:
	\begin{align*}
			\forall v_1 \ldots \forall v_k((&\phi(o, v_1,\ldots, v_k) \land {}\\
			&	(\forall x \in N)(\phi(x,v_1, \ldots, v_k) \lif \phi(s(x), v_1,\ldots, v_k))) \lif {}\\
			&\hspace{3em} (\forall x \in N)\phi(x, v_1,\ldots, v_k))
	\end{align*}
Cetha yen ukara ``duwe paramèter'' ndadekake bab iki luwih gampang
diwaca. (Ing \olref[sth][][]{part}, piranti iki bakal kerep banget
digunakake.)

Bali menyang aljabar Dedekind: saka saben aljabar Dedekind, kita uga
bisa nemtokake fungsi aritmetis lumrah kanggo panggunggungan,
ping-pingan, lan eksponensiasi. Nanging iki ora prasaja lan mbutuhake
teknik \emph{definisi rekursif}. Teknik iku bakal ditepungake lan
dibenerake adoh mengko ing konteks kang luwih umum. (Pamaca kang
seneng bisa bali menyang iki sawise \olref[sth][ord-arithmetic][]{chap},
utawa maca cara liyane, kayata \citealt[kaca~95--8]{Potter2004}.) Nanging,
yen $N, s, o$ mbentuk aljabar Dedekind, pungkasane kita bakal bisa
netepake:
\begin{align*}
	{a} + {o} &= {a} & & & {a} \times {o} &= {o} & & & {a}^{o} &= s(o)\\
{a} + {s}({b}) &= {s}({a}+{b}) &&& {a} \times {s}({b}) &= ({a}\times {b}) + {a}  & & & {a}^{{s}({b})} &= {a}^{b} \times {a}
\end{align*}
lan nuduhake yen fungsi-fungsi iki tumindak kaya kang dikarepake.

\end{document}
